% Generated packaging derivative; source remains byte-preserved.



\paragraph{Evidence and scope.}
The statements below are proved on paper, with exact rational finite checks
in the accompanying script. They have no Lean formal counterpart and no
external peer review. No novelty claim is made. They concern a restricted
class of population procedures and do not resolve unrestricted minimax risk.

\paragraph{Experiment.}
Let $\mu=\sum_{i=1}^k\pi_i\delta_{d_i}$, where $\pi_i>0$,
$\sum_i\pi_i=1$, and $0<d_i\le1/2$ are known. A new coin independently draws
$D\sim\mu$ and $S\mid D=d_i\sim\operatorname{Bernoulli}(q_i)$, then supplies
conditionally iid flips with success probability $1/2+(2S-1)D$.
The gap and sign are hidden. For known $0<\rho\le1/4$, the parameter space is
the compact convex set
\[
 Q=\{q\in[0,1]^k:\ \theta(q):=\textstyle\sum_i\pi_iq_i\le\rho\}.
\]
All coin observations, including pilot observations, are charged. Arithmetic
and internal randomization are free. The three-flip construction below has a
deterministic work count; no random-work approximation is used.

\begin{theorem}[A saddle point for a fixed scalar control variate]\label{thm:saddle}
Fix a single-coin observation rule and real features $Y,Z$ of its observations
and internal randomization. Suppose they have finite second moments under
each conditional law $(D,S)=(d_i,s)$, and for all $q\in Q$,
\[
 \E_qY=\theta(q),\qquad \E_qZ=0,\qquad \E_qZ^2=v>0,
\]
where $v$ does not depend on $q$. Put
$f(q,\beta)=\Var_q(Y-\beta Z)$ and
$F(\beta)=\max_{q\in Q}f(q,\beta)$.
There exist $\beta_*\in\mathbb R$ and $q_*\in Q$ such that
\[
 f(q_*,\beta)\ \ge\ f(q_*,\beta_*)=F(\beta_*)=\min_{b\in\mathbb R}F(b)
 \quad\text{for every }\beta\in\mathbb R.
\]
In particular,
\[
 \sup_{q\in Q}\inf_{\beta\in\mathbb R} f(q,\beta)
 =\inf_{\beta\in\mathbb R}\sup_{q\in Q} f(q,\beta).
\]
\end{theorem}

\begin{proof}
Write $c(q)=\E_q[YZ]$ and $V(q)=\Var_qY$. Conditional mixture probabilities
are affine functions of $q$. Therefore $c$ is affine and continuous,
$V(q)=\E_qY^2-\theta(q)^2$ is continuous and concave, and
\[
 f(q,\beta)=V(q)-2\beta c(q)+v\beta^2
\]
is concave in $q$ and strictly convex in $\beta$.
On compact $Q$, let $C=\max_q|c(q)|<\infty$. Since $V(q)\ge0$,
$F(\beta)\ge v\beta^2-2C|\beta|$, so $F$ is coercive. It is continuous,
being a maximum over a compact set of functions whose joint formula is
continuous. It has a global minimizer $\beta_*$.

Let $\mathcal A=\{q:f(q,\beta_*)=F(\beta_*)\}$ be its nonempty compact active
set. We show that the active slopes
$g(q)=2v\beta_*-2c(q)$ include one nonpositive and one nonnegative slope.
If every active slope were positive, compactness would give a neighborhood
$U$ of $\mathcal A$ on which $g\ge\eta>0$. On the compact complement
$Q\setminus U$, the value at $\beta_*$ is at most $F(\beta_*)-\gamma$ for
some $\gamma>0$ (if the complement is empty, omit it). For sufficiently small
$t>0$, the identity
\[
 f(q,\beta_*-t)=f(q,\beta_*)-t g(q)+vt^2
\]
then makes every value strictly smaller than $F(\beta_*)$: on $U$ use
$t\eta-vt^2>0$, and on its complement use the uniform bound on $|g|$ and
the gap $\gamma$. This contradicts minimality. The all-negative case follows
by increasing $\beta_*$. Thus there are $q_1,q_2\in\mathcal A$ and
$\lambda\in[0,1]$ for which
$\lambda g(q_1)+(1-\lambda)g(q_2)=0$.

Set $q_*=\lambda q_1+(1-\lambda)q_2\in Q$. Concavity gives
\[
 f(q_*,\beta_*)\ge\lambda f(q_1,\beta_*)+(1-\lambda)f(q_2,\beta_*)
 =F(\beta_*).
\]
The opposite inequality is the definition of $F$. Moreover $g(q_*)=0$
by affinity of $c$, so $\beta_*$ minimizes the quadratic $f(q_*,\cdot)$.
This proves the saddle assertion and both directions of the equality.
\end{proof}

\begin{theorem}[Independent pilot learning does not improve this minimax risk]
Under Theorem~\ref{thm:saddle}, take $n\ge1$ independent fresh production
coins, obtaining iid pairs $(Y_i,Z_i)$. A pilot may use arbitrary additional
coins and observations to choose an almost surely finite real coefficient
$\widehat\beta$, but its complete transcript and randomization must be
independent of the production coins under each $q$. Let
\[
 \widehat\theta=\frac1n\sum_{i=1}^n(Y_i-\widehat\beta Z_i),\qquad
 F_*:=F(\beta_*).
\]
Then, allowing infinite mean squared error,
\[
 \sup_{q\in Q}\E_q(\widehat\theta-\theta(q))^2\ge F_*/n.
\]
The fixed coefficient $\beta_*$ attains equality without any pilot. If the
features use exactly $\ell$ flips per coin and a pilot spends $\ell m$ flips,
then a fixed-coefficient estimate from $m+n$ coins instead has worst-case
MSE $F_*/(m+n)$ at the same deterministic total work.
\end{theorem}

\begin{proof}
For each realized finite coefficient, the production estimate has conditional
mean $\theta(q)$ and conditional squared error $f(q,\widehat\beta)/n$.
Taking expectations of this nonnegative conditional second moment is valid
even when it is infinite. At the specific $q_*$ in
Theorem~\ref{thm:saddle}, every possible coefficient has
$f(q_*,\widehat\beta)\ge F_*$. A fixed $\beta_*$ has worst-case variance
$F_*$ per production coin and is unbiased. The work comparison follows by
counting coins. The last assertion requires that the stated pilot work is
indeed deterministic; no conversion from expected work is implicit.
\end{proof}

\paragraph{Excluded procedures.}
The final estimate here is the raw affine average. Clipping it, changing the
production feature or sampling rule after seeing a pilot, adaptively selecting
the production sample size, accepting bias, and nonlinear aggregation are all
outside the theorem. Thus this result does not transfer a single-coin lower
bound to unrestricted population procedures.

\begin{proposition}[Explicit two-gap three-flip construction]\label{prop:two}
Suppose $k=2$, $0<d_1<d_2\le1/2$, and put $x_i=2d_i$. Take exactly three
flips on a fresh coin and put $R_j=2X_j-1$,
\[
 A=(R_1+R_2+R_3)/3,\quad B=R_1R_2R_3,\quad
 C=(R_1R_2+R_1R_3+R_2R_3)/3.
\]
Define
\[
 a=\frac{x_1^2+x_1x_2+x_2^2}{2x_1x_2(x_1+x_2)},\qquad
 b=-\frac1{2x_1x_2(x_1+x_2)},\qquad
 m_2=\pi_1x_1^2+\pi_2x_2^2,
\]
and $Y=1/2+aA+bB$, $Z=C-m_2$. Then $\E_qY=\theta(q)$ and
$\E_qZ=0$ for every $q$. Moreover
\[
 v:=\E_q Z^2=(1-m_2)(m_2+1/3)>0.
\]
Write $m_j=\sum_i\pi_i(2q_i-1)x_i^j$ for odd $j$, and
\[
 K=a^2(1+2m_2)/3+2abm_2+b^2.
\]
The exact moments are
\begin{align*}
 V(q)&=K-1/4+\theta(q)-\theta(q)^2,\\
 c(q)&=a\{(2/3-m_2)m_1+m_3/3\}+b(m_1-m_2m_3).
\end{align*}
Consequently these features satisfy Theorem~\ref{thm:saddle}.
\end{proposition}

\begin{proof}
Conditionally on $D=d_i,S=s$, the $R_j$ are iid with mean
$x=(2s-1)x_i$. Hence $\E A=x$, $\E B=x^3$, $\E C=x^2$.
The displayed coefficients solve $ax_i+bx_i^3=1/2$ for both $i$, so
$\E[Y\mid D=d_i,S=s]=s$.
The identity $\E[C\mid D=d_i,S=s]=x_i^2$ proves the assertion for $Z$.
Direct multiplication using $R_j^2=1$ gives
\[
 \E A^2=\E C^2=(1+2m_2)/3,\quad \E B^2=1,\quad
 \E AB=m_2,
\]
and
\[
 \E AC=2m_1/3+m_3/3,\qquad \E BC=m_1.
\]
These yield the claimed variance and covariance formulas. Since the two
gaps are distinct and both masses are positive, $0<m_2<1$, even when
$d_2=1/2$. Thus $v>0$. All features are bounded for the fixed positive gaps,
and no sampled latent value is used by the computation.
\end{proof}

\begin{proposition}[A fully specified pilot implementation]
Use $m\ge1$ independent pilot coins with the features in
Proposition~\ref{prop:two}. Set
\[
 \widehat\beta=\frac1{mv}\sum_{j=1}^mY_jZ_j.
\]
On $n\ge1$ independent production coins return the raw average of
$Y-\widehat\beta Z$. It is unbiased for all $q$, terminates after exactly
$3(m+n)$ flips, and has exact MSE
\[
 \frac1n\left(V(q)-\frac{c(q)^2}{v}
             +\frac{\Var_q(YZ)}{mv}\right).
\]
\end{proposition}

\begin{proof}
The pilot is independent of production, with
$\E\widehat\beta=c(q)/v$ and
$\Var(\widehat\beta)=\Var_q(YZ)/(mv^2)$. Completing the square in
$f(q,\beta)=V(q)-c(q)^2/v+v(\beta-c(q)/v)^2$ and then averaging proves
the formula. Conditional unbiasedness proves unconditional unbiasedness.
All moments are finite for fixed gaps, and the work count is deterministic.
\end{proof}

\paragraph{Interpretation.}
The explicit estimator is an exact hidden-gap construction, not an efficient
rare-event algorithm: its interpolation coefficients grow as the gaps shrink,
and it can have a large variance even at $\theta=0$. The no-gain theorem
shows why merely learning a coefficient for fixed unbiased features cannot
be the sought minimax improvement. It leaves acquisition adaptation and
nonlinear estimators as separate questions. No claim is made that the
control-variate idea or the convex saddle argument is new.
